% !PRINTED 384
\weberpage{384}
%% [running head, not body text -- commented out 2026-08-07] Page furniture (running
%% head, folio, signature) is excluded from this edition by standing policy.
%% 384 \hfill Fifteenth Chapter. \hfill § 104.

%<<CONT
quadratic divisor of $m$, which one again decomposes into two coprime
quadratic factors, of which the one is taken to $n$, the other
to $n'$, and thereby $(\delta, n)$ and
$(\delta', n')$ are not altered.

Accordingly one can determine a character $\chi(\delta, n)$ of an arbitrary
number coprime to $Q$ in the following way.

One sets:

\weq{18}{n = \varepsilon n'}

and lets $\varepsilon$ denote the product of the powers dividing $n$ of
those primes which at the same time divide $\delta$; then $\varepsilon$ is relatively
prime to $\delta'$ and $n'$ relatively prime to $\delta$, and one sets:

\weq{19}{\chi(\delta, n) = (\delta', \varepsilon)(\delta, n'),}

by which $\chi(\delta, n)$ is uniquely determined and always different from zero. If then $(a, b, c)$ is a representative of the class $A$,
then

\weq{20}{\chi(\delta, ax^2 + bxy + cy^2) = \chi(\delta, A),}

in which $x, y$ are arbitrary integers with or without a common
divisor, for which $ax^2 + bxy + cy^2$ becomes coprime to $Q$.

Concerning the symbol $\chi(\delta, n)$ so defined there holds a series of
theorems which we have still to derive.

6. \gesperrt{If $n_1$ and $n_2$ are any two numbers coprime to $Q$,
then}

\weq{21}{\chi(\delta, n_1)\chi(\delta, n_2) = \chi(\delta, n_1n_2).}

For if we set, according to (18),
\[
n_1 = \varepsilon_1 n_1', \quad n_2 = \varepsilon_2 n_2',
\]
then
\[
n_1 n_2 = \varepsilon_1 \varepsilon_2 n_1' n_2',
\]
and this passes into (18) if one sets $n = n_1n_2$, $\varepsilon = \varepsilon_1\varepsilon_2$,
$n' = n_1'n_2'$. Thus (21) follows from § 103 (8).

7. \gesperrt{If one lets $e$ run through the sequence of divisors of $n$, and sets $n = ee'$, then}

\weq{22}{\chi(\delta, n)\overset{e}{\Sigma}(D, e) = \overset{e}{\Sigma}(\delta, e)(\delta', e').}

If formula (22) holds for $n = n_1$ and $n = n_2$, then it holds,
under the assumption that $n_1, n_2$ are relatively prime, also
for $n = n_1n_2$ [by (21)]; it therefore needs to be proved only for a
prime power $n = p^k$, when $p$ does not divide $Q$
%CONT>>
